\documentclass[11pt]{article}
\usepackage[margin=1in]{geometry}
\usepackage{amsmath,amssymb}
\usepackage{booktabs}
\usepackage{array}
\usepackage{float}
\usepackage{xcolor}
\usepackage[colorlinks=true,linkcolor=blue,urlcolor=blue]{hyperref}
\usepackage{enumitem}

\newcommand{\bigO}{\mathcal{O}}
\newcommand{\states}{\lvert S\rvert}
\newcommand{\code}[1]{\texttt{#1}}

\title{\textbf{Complexity analysis of the RTIGER core optimizations}\\[2pt]
\large A per-step big-$\mathcal{O}$ accounting of the \code{optimize-julia-core} fork}
\author{Gemma Lab (\code{faustovrz/RTIGER})}
\date{\today}

\begin{document}
\maketitle

\begin{abstract}
\noindent
This note derives, step by step, the asymptotic cost of the rigid-HMM (rHMM)
Baum--Welch fit that RTIGER performs, and the big-$\bigO$ improvement contributed
by each optimization in the \code{optimize-julia-core} fork. Equation numbers
refer to the supplementary text of Campos-Martin \emph{et al.}\ (2023).%
\footnote{Campos-Martin R, Schmickler S, Goel M, Schneeberger K, Tresch A.
\emph{Reliable genotyping of recombinant genomes using a robust hidden Markov
model.} Plant Physiology 2023;192(2):821--836.
\texttt{doi:10.1093/plphys/kiad191}; PMC10231367.
\url{https://doi.org/10.1093/plphys/kiad191}} Every optimization preserves the
arithmetic of the model (Sections~2--4 of the supplement); only the cost of
evaluating it changes. Two of the optimizations remove an asymptotic factor
(the rigidity window $r$, and the marker count $T$ in the emission line search),
one removes a memory factor (the sample count $N$), and two are
constant-factor (allocation) wins that leave the complexity class unchanged.
\end{abstract}

\section{Notation and cost model}

\begin{table}[h]
\centering
\begin{tabular}{c l}
\toprule
symbol & meaning \\
\midrule
$T$        & markers (positions) in one observation chain \\
$C$        & number of chains (samples $\times$ chromosomes); $T_{\text{tot}}=\sum_c T_c$ \\
$N$        & number of samples (so $C$ and $T_{\text{tot}}$ grow $\propto N$) \\
$\states=s$ & number of hidden states ($=3$: paternal / heterozygous / maternal) \\
$r$        & rigidity; augmented rigid state space $S'=S\times\{1,\dots,r\}$, $\lvert S'\rvert = s\,r$ \\
$D$        & number of \emph{distinct} emission pairs $(k,n)$ in a chain \\
$E$        & objective evaluations per state in one emission-M-step line search \\
\bottomrule
\end{tabular}
\caption{Cost parameters. For low-coverage genotyping data $D\ll T$: the read
depth $n$ is small, so the distinct $(k,n)$ pairs number a few dozen regardless
of $T$.}
\end{table}

\noindent
The emission is a Beta--Binomial (supplement \S3.3),
$\psi_s(k;n,a_s,b_s)=\binom{n}{k}\,\mathrm{B}(k+a_s,n-k+b_s)/\mathrm{B}(a_s,b_s)$,
whose evaluation costs a handful of \code{lgamma} calls --- the dominant scalar
op of the fit. A full fit is
\[
\text{cost} \;=\; (\#\text{EM iterations})\times\sum_{c=1}^{C}\big(\text{E-step}_c + \text{M-step}_c\big),
\]
so we analyse the per-iteration, per-chain cost and then note how it aggregates.

\section{The auxiliary window product $\Psi^t_k$ (\code{productpsi})}

The rigid recursions repeatedly need the windowed emission product
(supplement Eq.~21)
\begin{equation}
\Psi^{t}_k \;=\; \prod_{u=\max(t-r+1,1)}^{t}\psi_k(o_u),
\qquad 1\le t\le T .
\tag{21}
\end{equation}
Computing each of the $T\cdot s$ values by re-multiplying its length-$r$ window
costs $\bigO(T\,s\,r)$. In log space $\log\Psi^t_k=\sum_{u}\log\psi_k(o_u)$ is a
fixed-width sliding sum: maintain a running window, add the entering term and
drop the leaving term in $\bigO(1)$ per step.
\[
\boxed{\;\bigO(T\,s\,r)\;\longrightarrow\;\bigO(T\,s)\;}
\qquad\text{(removes the rigidity factor }r\text{; commit \code{d725651}).}
\]
At the production rigidity $r=250$ this is a $\sim r$-fold reduction of the
$\Psi$ term.

\section{Emission log-pdf memoization (\code{getlogpsi})}

Each forward, backward and Viterbi step evaluates $\psi_s(o_t)$ at every position
(Eqs.~22--27, 75). Done literally, that is $\bigO(T\,s)$ Beta--Binomial
evaluations per iteration --- each one several \code{lgamma} calls. But
$\psi_s(o_t)$ depends on $o_t=(k_t,n_t)$ only, and there are only $D$ distinct
pairs. Memoizing $\log\psi_s$ over the distinct $(k,n)$ replaces the expensive
evaluations with cheap table lookups:
\[
\boxed{\;\bigO(T\,s)\ \text{Beta--Binomial evals}\;\longrightarrow\;\bigO(D\,s)\ \text{evals} + \bigO(T\,s)\ \text{lookups}\;}
\]
i.e.\ the costly part shrinks by the factor $T/D$ (commit \code{d725651}).

\section{Emission M-step line search (the dominant cost)}

The Beta--Binomial parameters are fit by a 1-D line search over $\tau_s$ that
maximizes (supplement Eqs.~38, 43)
\begin{equation}
Q(\tau_s)\;\propto\;\sum_{c=1}^{C}\sum_{t=1}^{T_c}\gamma^{t,c}_s\,
\log\psi_s\!\big(k^{c}_t;n^{c}_t,\,a_s=m_s\tau_s,\,b_s=(1-m_s)\tau_s\big).
\tag{43}
\end{equation}
\textbf{Original.} Each of the $E$ line-search evaluations rebuilds a
Beta--Binomial object and sums $\log\psi_s$ over all $T_{\text{tot}}$ markers:
$\bigO(E\,s\,T_{\text{tot}})$ Beta--Binomial constructions per iteration. The
supplement notes the forward/backward/Viterbi kernels are the ``time-critical''
parts; in practice this repeated full-marker objective is what dominates ---
empirically $\sim 99.97\%$ of upstream runtime.

\textbf{Optimized.} Because the sum only depends on the data through the
$\gamma$-weighted counts at each distinct pair, regroup it once:
\begin{equation}
Q(\tau_s)=\sum_{(k,n)\in\mathcal D}\!\underbrace{\Big(\!\!\sum_{t:\,(k_t,n_t)=(k,n)}\!\!\gamma^{t}_s\Big)}_{G_s(k,n)\ \text{(pre-summed once)}}\;\log\psi_s(k;n,a_s,b_s).
\end{equation}
The weights $G_s(k,n)$ are accumulated in one $\bigO(s\,T_{\text{tot}})$ pass;
each line-search evaluation then costs $\bigO(s\,D)$:
\[
\boxed{\;\bigO(E\,s\,T_{\text{tot}})\;\longrightarrow\;\bigO\big(s\,T_{\text{tot}}+E\,s\,D\big)\;}
\]
Since $E\,D\ll T_{\text{tot}}$, the M-step collapses from ``$E$ passes over all
markers, rebuilding objects'' to ``one binning pass'' --- a measured
$\sim$\,1500$\times$ at production scale (commit \code{16b8a65}). After this, the
M-step is no longer the bottleneck.

\section{Forward / backward and Viterbi (constant-factor)}

The rigid forward (Eqs.~22--23), backward (Eqs.~26--27) and Viterbi (Eq.~75)
\[
\varphi^{t}_k=\max_{j\in S}
\begin{cases}\psi_k(o_t)\,a_{kk}\,\varphi^{t-1}_k & j=k\\[2pt]
\Psi^{t}_k\,a_{jk}\,\varphi^{t-r}_j & j\ne k\end{cases}
\tag{75}
\]
each take a $\max$/sum over $s$ predecessors at every position, giving
$\bigO(T\,s^2)$ per chain (the dwell-counter shifts for $k<r$ are $\bigO(1)$
copies once $\Psi$ is available). The optimizations here --- in-place scalar
log-sum-exp (Eq.~80) and an allocation-free in-place arg-max --- do \emph{not}
change the class:
\[
\bigO(T\,s^2)\;\longrightarrow\;\bigO(T\,s^2),
\qquad\text{constant-factor: }\sim 5\times\ (\text{fwd/bwd, commit \code{44ed85b}}),\ \sim 12\times\ (\text{Viterbi, \code{6cf85ca}}),
\]
with $\sim$90$\times$ fewer allocations in the forward/backward kernel. Once the
emission M-step is fixed (\S4), these $\bigO(T\,s^2)$ kernels are what the
per-iteration cost reduces to --- hence \emph{linear in $T$}.

\section{Memory: streaming the M-step (\code{eb79933})}

The M-step needs the data only through pooled sufficient statistics that are
\emph{sums} over chains (Eqs.~40--42): $\sum_c\zeta$, $\sum_c\gamma^{1}$, and the
$\gamma$-weighted counts $G_s(k,n)$.

\textbf{Original.} \code{EM()} held, for \emph{all} $C$ chains at once, the
per-chain arrays $\zeta$ ($T\!\times\!s^2$), $\gamma$ ($s\!\times\!T$) and
$\alpha,\beta,\psi$ ($s r\!\times\!T$), pooling only at the end:
\[
\text{peak} = \bigO\!\Big(\textstyle\sum_c T_c\,(s^2+s r)\Big)=\bigO\big(N\,T\,(s^2+s r)\big)\quad\text{(linear in $N$).}
\]
\textbf{Streaming.} Fold each chain's contribution into the running sums inside
the E-step loop and discard its arrays immediately (the public
\code{@Probabilities} slot, which nothing reads, is no longer filled):
\[
\boxed{\;\bigO\big(N\,T\,(s^2+s r)\big)\;\longrightarrow\;\bigO\big(T_{\max}\,(s^2+s r)\big)+\bigO(T_{\text{tot}})\;}
\]
peak scratch becomes one chain's worth, \emph{independent of $N$}; the only
$N$-linear residual is the input count matrices that must be held. Measured:
$\sim$33\,GB projected $\to\sim$3.6\,GB at $1400\times50$k.

\section{Summary}

The per-step time and memory complexity is consolidated in the table below.

\begin{table}[H]
\centering
\renewcommand{\arraystretch}{1.3}
\begin{tabular}{l l l c}
\toprule
optimization & before (per iter) & after (per iter) & gain \\
\midrule
$\Psi$ window (\code{productpsi})        & $\bigO(T s r)$            & $\bigO(T s)$                       & drop $r$ \\
emission log-pdf (\code{getlogpsi})      & $\bigO(T s)$ evals        & $\bigO(D s)$ evals                 & $T/D$ \\
emission M-step (\code{16b8a65})         & $\bigO(E s T_{\text{tot}})$ & $\bigO(s T_{\text{tot}}+E s D)$  & $\sim$1500$\times$ \\
forward/backward (\code{44ed85b})        & $\bigO(T s^2)$            & $\bigO(T s^2)$                     & $\sim$5$\times$ (const.) \\
Viterbi (\code{6cf85ca})                 & $\bigO(T s^2)$            & $\bigO(T s^2)$                     & $\sim$12$\times$ (const.) \\
\midrule
memory, streaming (\code{eb79933})       & $\bigO(N T (s^2{+}s r))$  & $\bigO(T_{\max}(s^2{+}s r))$       & drop $N$ \\
\bottomrule
\end{tabular}
\caption{Per-EM-iteration time complexity (top) and peak-memory complexity
(bottom). $s=3$ is constant; the asymptotically meaningful factors are $T$, $r$,
$D$, $E$ and $N$.}
\end{table}

\section{Reconciliation with the measured scaling}

A marker-scaling head-to-head on one shared 109{,}703-locus panel (decimated by
odd index to five sizes, both cores run \emph{to convergence}: \texttt{eps}$=0.01$,
$r=2$, identical deterministic init) gives empirical per-iteration-mean exponents
\[
\text{original}\approx \text{markers}^{2.20},
\qquad
\text{optimized}\approx \text{markers}^{1.34}.
\]
These match the analysis:
\begin{itemize}[nosep]
\item \textbf{Optimized.} With the M-step reduced to a binning pass plus an
$\bigO(E s D)$ search ($D$ nearly flat), per-iteration cost is dominated by the
$\bigO(T s^2)$ forward/backward/Viterbi kernels --- essentially \emph{linear in
$T$}.
\item \textbf{Original.} Per-iteration cost is dominated by the emission line
search $\bigO(E s T_{\text{tot}})$ with a Beta--Binomial rebuilt per marker, and
$E$ (the \code{Optim} evaluation count) is data-dependent and grows with scale.
The result is an empirically super-linear, erratic per-iteration cost
($\sim$markers$^{2.2}$, with iteration-to-iteration spikes) rather than a clean
$\bigO(T)$ --- exactly what the grouping in \S4 removes.
\end{itemize}
So per-iteration cost goes from empirically $\sim$quadratic and erratic to
provably linear ($\bigO(T s^2)$), and the full-fit speed-up \emph{widens} with
problem size: \textbf{$\sim$30$\times$} at $6{,}857$ markers/sample to
\textbf{$\sim$405$\times$} at $109{,}703$ ($3.4$\,h\,$\to$\,$30$\,s, both in
$19$ EM iterations).

\paragraph{Equivalence.} Crucially the speed-up changes nothing in the result.
Run to convergence at all five sizes, the two cores take the \emph{same} number
of EM iterations ($5,7,11,16,19$ for opt and orig alike), return \emph{identical}
Viterbi paths ($329{,}109/329{,}109$ positions at the full panel; $0$ mismatches
at every size), and fitted parameters agree to $\le 3\times10^{-6}$ --- float
summation-order, not an algorithmic difference. The optimized core is the same
estimator, computed at $\bigO(T)$ instead of empirically $\bigO(T^2)$ in markers.

\section{Parallelism: the E-step across chains (a constant, not an exponent)}

The E-step is independent across chains (samples $\times$ chromosomes) — the
M-step only pools their summed sufficient statistics — so it is embarrassingly
parallel. With $P$ worker threads the per-iteration wall time is
\[
\bigO(T s^2)\;\longrightarrow\;\bigO\!\left(\tfrac{T s^2}{P}\right),
\qquad P \le \min(\text{Julia pool},\ \text{physical cores}),
\]
i.e.\ the \emph{constant} drops by $P$ but the \emph{exponent} is unchanged: on a
log--log plot the optimized line shifts down by $\log P$ with the same slope.
Parallelism cannot beat $\bigO(N)$ — dividing a fixed amount of work among a
\emph{bounded} number of workers is a constant-factor speed-up, not a
complexity-class change. (A sub-linear exponent would need $P$ to grow with the
data; the chromosome is the natural, fixed parallel grain, so it does not.)

In practice the realized speed-up is well below $P$ because the kernel is
\textbf{memory-bandwidth bound} — it streams the large $\alpha,\beta,\zeta$ arrays
with little arithmetic per element, so a few cores saturate the shared bandwidth.
Measured on a 4-performance-core Apple M4: $\sim$1.4$\times$ at $P=4$ on both real
\emph{Arabidopsis} (15 chains) and maize (10 chains) data, stable across chain
counts, with GC $\sim$2.5\% (not the limiter) and $P=8$ no better than $P=4$.
Determinism is preserved by a fixed-order reduce and a convergence guard band;
\code{threads}=1 (default) is the bit-identical serial path and \code{threads}$>$1
is Viterbi-identical. Commits \code{664eb33}, \code{be437d1}.

\end{document}
